%! Author = Maximilian Fülle
%! Date = 29.05.2026


\chapter{Models of curves over valued fields}\label{ch:models-of-curves-over-valued-fields}

This chapter introduces models of smooth projective curves over valued fields, together with their morphisms and the resulting partial order by domination.
After discussing normality and basic examples, we attach a geometric valuation and a multiplicity to every irreducible component of the special fibre; the resulting finite set of valuations is the starting point for the correspondence developed in Chapter~\ref{ch:correspondence-between-models-and-valuations}.

\section{Valued fields, curves, and models}\label{sec:valued-fields-curves-and-models}
We begin by fixing the valued-field notation and standing hypotheses used throughout Part~\ref{part:models-and-valuations}, and then introduce the basic notion of a model of a curve.

Let $K$ be a field endowed with an additive valuation of rank $1$, $v_K: K \to \hat{\RR} \coloneqq \RR\cup \lbrace\infty\rbrace$.
Let $\Gamma_K \coloneqq v_K(K^*) \subset \RR$ be the value group of $v_K$, $\OK \subset K$ its valuation ring with
maximal ideal $\mK$ and $k \coloneqq \OK /\mK$ its residue field.

For most of the first chapter, we make the following additional assumptions:
\begin{assumption} \label{ass:v_K}
    \begin{enumerate}[label=\textup{(\alph*)}]
        \item $v_K$ is discrete, i.e.\ $\Gamma_K$ is a discrete subgroup of $\RR$,
        \item $K$ is complete with respect to $v_K$,
        \item $k$ is perfect.
    \end{enumerate}
\end{assumption}

Throughout the monograph, $X$ will denote a \textit{nice} curve over $K$,
that is a smooth, projective and absolutely irreducible curve over $K$.

\begin{definition} \label{def:model}
    A \textit{model} $\X$ of $X$ is a proper, flat and normal $\OK $-scheme
    with an identification of its generic fibre $\X_\eta$ with $X$.
\end{definition}

\begin{example}
    Let $K = \QQ_p$, $\OK  = \ZZ_p$ and $X = \PP^1_K$.
    Then $\X = \PP^1_{\OK}$ is a model of $X$ with the evident identification
    $\X_\eta = \X_K \cong \PP^1_K$.
    The special fibre $\X_s$ has one irreducible component which is isomorphic to $\PP^1_k$.
\end{example}

\begin{remark} \label{rem:flatIffIntegral}
    The condition of a model being flat is equivalent to it being integral:
    if $\X$ is integral and its structural map is compatible with the given generic fibre, then every local ring $\OO_{\X,x}$ is a domain containing $\OK$.
    Thus $\OO_{\X,x}$ is a torsion-free module over the discrete valuation ring $\OK$, hence flat by~\cite[Corollary 1.2.14]{Liu}.
    Conversely, the integrality of a model $\X$ of $X$ follows from its flatness and the integrality of $X$ by~\cite[Proposition 4.3.8]{Liu}.
\end{remark}


\section{The category of models and modifications}\label{sec:category-of-models-and-modifications}
We now introduce the morphisms used to compare models of a fixed curve.
They turn the collection of models into a partially ordered system in which any two models admit a common dominating model.

\begin{definition}\label{def:modification-of-models}
    Let $\X$ and $\X'$ be models of $X$.
    A \textit{modification} $\varphi\colon\X'\to\X$ is a morphism of $\OK$-schemes whose restriction to the generic fibres is the identity on $X$ under the fixed identifications $\X'_\eta\cong X\cong\X_\eta$.
    In this situation, we say that $\X'$ \textit{dominates} $\X$.
\end{definition}

\begin{lemma}\label{lem:uniqueness-of-modifications}
    Between two fixed models of $X$, there is at most one modification.
\end{lemma}

\begin{proof}
    Let $\varphi,\psi\colon\X'\to\X$ be modifications.
    They agree on the generic fibre $\X'_\eta$, which is dense in the integral scheme $\X'$.
    Since $\X$ is separated, the equalizer of $\varphi$ and $\psi$ is closed in $\X'$.
    It contains the dense generic fibre and is therefore all of $\X'$, so $\varphi=\psi$.
\end{proof}

\begin{proposition}\label{prop:elementary-properties-of-modifications}
    Let $\varphi\colon\X'\to\X$ be a modification of models.
    Then $\varphi$ is proper, dominant, and surjective.
    Moreover, there is a finite set $S\subset\X_s$ of closed points such that $\varphi$ is an isomorphism over $\X\setminus S$ and, for every $x\in S$, the fibre $\varphi^{-1}(x)$ is a connected projective curve over $\kappa(x)$.
\end{proposition}

\begin{proof}
    Let $\varphi\colon\X'\to\X$ be a modification.
    Since $\X$ is separated over $\Spec\OK$, the graph morphism
    \[
        \Gamma_\varphi\colon\X'\longrightarrow\X'\times_{\Spec\OK}\X
    \]
    is a closed immersion.
    The projection $\X'\times_{\Spec\OK}\X\to\X$ is proper because it is the base change of the proper morphism $\X'\to\Spec\OK$.
    Their composite is $\varphi$, so $\varphi$ is proper.
    Its image contains the generic fibre $X$, which is dense in $\X$, and hence $\varphi$ is dominant.
    Finally, a proper morphism has closed image, so the image of $\varphi$ is both dense and closed in $\X$ and must equal $\X$.

    A proper birational morphism to a normal scheme is an isomorphism at every codimension-one point of the target; this follows by applying the valuative criterion of properness to the corresponding discrete valuation ring.
    Hence the complement of the maximal open subset of $\X$ over which $\varphi$ is an isomorphism is a closed subset of codimension at least two.
    Since $\X$ is a two-dimensional Noetherian scheme and $\varphi$ is already an isomorphism on the generic fibre, this complement is a finite set $S\subset\X_s$ of closed points.

    By the connectedness form of Zariski's Main Theorem, the fibres of $\varphi$ are connected; see~\cite[\href{https://stacks.math.columbia.edu/tag/0AY8}{Lemma 0AY8}]{Stacks}.
    If $x\in S$, the fibre over $x$ cannot be zero-dimensional: otherwise $\varphi$ would be finite in a neighbourhood of $x$, and a finite birational morphism to the normal scheme $\X$ would be an isomorphism there.
    Thus $\varphi^{-1}(x)$ is a proper one-dimensional scheme over $\kappa(x)$, hence a projective curve.
\end{proof}

We write
\[
    \X\leq\X'
\]
if there is a modification $\X'\to\X$; thus the larger model is the one that dominates the smaller model.
Here and below, models that are isomorphic as models of $X$ are identified when discussing this relation.

\begin{proposition}\label{prop:partial-order-of-models}
    The relation $\leq$ is a partial order on the set of isomorphism classes of models of $X$.
    The category of models and modifications is equivalent to the partially ordered set regarded as a category.
\end{proposition}

\begin{proof}
    The identity morphism gives reflexivity, and composition of modifications gives transitivity.
    If $\X\leq\X'$ and $\X'\leq\X$, there are modifications $\varphi\colon\X'\to\X$ and $\psi\colon\X\to\X'$.
    Both composites are modifications and hence, by Lemma~\ref{lem:uniqueness-of-modifications}, are the respective identity morphisms.
    Thus $\varphi$ and $\psi$ are mutually inverse, which proves antisymmetry on isomorphism classes.
    The same uniqueness lemma says that there is at most one morphism between any two models, so choosing one representative of each isomorphism class gives the asserted equivalence of categories.
\end{proof}

\begin{proposition}\label{prop:common-dominating-model}
    Any two models $\X_1$ and $\X_2$ of $X$ admit a common dominating model $\X_3$.
\end{proposition}

\begin{proof}
    Consider the diagonal morphism determined by the fixed identifications of the generic fibres,
    \[
        X\longrightarrow\X_1\times_{\Spec\OK}\X_2.
    \]
    Let $\mathcal Z$ be the schematic closure of its image.
    The product $\X_1\times_{\Spec\OK}\X_2$ is proper over $\Spec\OK$, and hence so is its closed subscheme $\mathcal Z$.
    Moreover, $\mathcal Z$ is integral and its generic fibre is the diagonal copy of $X$.
    It is therefore flat over the discrete valuation ring $\OK$ by Remark~\ref{rem:flatIffIntegral}.

    Let $\X_3$ be the normalization of $\mathcal Z$ in $F_X$.
    The completeness assumption in Assumption~\ref{ass:v_K} enters here: it implies that $\OK$ is excellent, so the normalization is finite; see~\cite[Chapter 8, in particular Theorem 8.2.39 (d)]{Liu}.
    Consequently, $\X_3$ is proper over $\Spec\OK$.
    It is integral and hence flat over $\OK$, and it is normal by construction.
    Since $X$ is normal, normalization does not change the generic fibre, so $(\X_3)_\eta=X$.
    Thus $\X_3$ is a model of $X$.
    Composing the normalization map with the two projections gives modifications $\X_3\to\X_i$ for $i=1,2$.
\end{proof}

\begin{remark}\label{rem:models-form-cofiltered-category}
    Proposition~\ref{prop:common-dominating-model} says that the partially ordered set of models is directed upward.
    Since modifications point from a dominating model to a dominated model, the category of models is correspondingly cofiltered and may be viewed as an inverse system.
    Common dominating models may also be produced through suitable blowups followed by normalization.
\end{remark}

Chapter~\ref{ch:correspondence-between-models-and-valuations} will identify this order with inclusion between the corresponding finite sets of geometric valuations.


\section{Normality and examples of models}\label{sec:normality-and-examples-of-models}
To recognize when a proper flat scheme is a model, it remains to understand its normality. We first recall a criterion in terms of the special fibre and then apply it to several examples.

\begin{definition}
    A \textit{pre-model} of $X$ is a flat and proper $\OK$-scheme $\X$ with $\X_\eta = X$.
    That is, $\X$ satisfies all the conditions of a model except that it need not be normal.
\end{definition}

The following normality criterion is a reformulation of~\cite[Lemma 2.1]{ArzdorfWewers} in the present notation.

\begin{lemma}\label{lem:conditions-for-normality}
    Let $\X$ be a pre-model of $X$.
    Then $\X$ is normal (i.e.\ a model) if and only if the following two conditions hold.
    \begin{enumerate}[label=\textup{(\alph*)}]
        \item The special fibre $\X_s$ has no embedded points, i.e.\ no point $x\in \X_s$ \label{lem:conditions-for-normality-i}
            satisfying $\m_x = \Ann{R}{r}$ for some $r\in R \coloneqq \OO_{\X_s,x}$ (see~\cite[Definition 7.1.6]{Liu}).
        \item The local ring $\OO_{\X,\xi}$ is a discrete valuation ring for every generic point $\xi \in \X_s$. \label{lem:conditions-for-normality-ii}
    \end{enumerate}
\end{lemma}

\begin{proof}
    By Serre's criterion for normality~\cite[Theorem 8.2.23]{Liu},
    we need to check the following properties:
    \begin{itemize}
        \item[$(\mathrm{R}_1)$] $\X$ is regular at $x$ when $\codim{x} \leq 1$, and
        \item[$(\mathrm{S}_2)$] $\depth{\OO_{\X,x}} \geq 2$
    \end{itemize}
    for all $x\in \X$. \\
    The points on $\X$ with codimension $1$ are precisely the points on the generic fibre, which is regular by assumption,
    and the generic point on the special fibre.
    The codimension-one local rings under consideration are one-dimensional Noetherian local domains.
    Such a ring is regular if and only if it is a discrete valuation ring, so $(\mathrm{R}_1)$ is equivalent to
    \ref{lem:conditions-for-normality-ii}.
    By~\cite[Proposition 8.2.11]{Liu}, condition $(\mathrm{S}_2)$ is equivalent to the condition
    that $\depth{\OO_{\X_s, x}} \geq 1$ for all $x \in \X_s$.
    This condition always holds at a generic point, and at any other point $x$ it is equivalent to the condition that $x$ is not an embedded point.
\end{proof}

\begin{example}
    \begin{enumerate}[label=\textup{(\alph*)}]
        \item Consider a pre-model with special fibre
            \begin{align}
                \X_s \coloneqq \Spec{k[x,y]/(x^2,xy)},
            \end{align}
            which is topologically equivalent to $\AA^1_k$.
            The point corresponding to the maximal ideal $(x,y)$ is an embedded point.
        \item We consider the pre-model of a smooth curve with an embedded point on the special fibre described in \cite[Example \textrm{III}.9.8.4]{Hartshorne}.
            Let $C$ be the curve in $\PP^3_{\QQ_p} = \Proj{\QQ_p[X,Y,Z,W]}$, which is given by the ideal
            \begin{align}
                I = \left(p^2(x+1) - z^2, px(x+1) - yz, xz - py, y^2 - x^2(x+1) \right)
            \end{align}
            on the affine chart $W \neq 0$.
            $p$ is not a zero divisor in $\ZZ_p[x,y,z]/I$.
            Thus, the projective $\ZZ_p$-scheme $\mathcal{C}$ defined by $I$ is flat and hence a pre-model of $C$.
            On the affine chart $W \neq 0$, the special fibre $\mathcal{C}_s$ is given by the ideal
            \begin{align}
                I_p = (z^2, yz, xz, y^2 - x^2(x+1))
            \end{align}
            and has support equal to the nodal cubic curve $y^2 = x^2(x+1)$.
            At any point where $x \neq 0$, $\mathcal{C}_s$ is reduced.
            However, in the stalk at $(0,0,0)$, the element $z$ is a nonzero nilpotent because $z^2=0$.
            Thus, $(0,0,0)$ is an embedded point, and $\mathcal{C}$ is not a model of $C$ by Lemma \ref{lem:conditions-for-normality}.
    \end{enumerate}
\end{example}


\begin{comment}
    \begin{example}
    In this example, we consider models of curves in the plane $\PP^2_K$.
    Such a (smooth, projective and irreducible) curve $X$ is defined by a homogeneous irreducible polynomial $F(X,Y,Z) \in K[X,Y,Z]$.
    We show that if $F(X,Y,Z)$ is integral (that is all coefficients of $F$ are in $\OK$) and primitive, a model of $X$ is given by
    \begin{align}
        \X = \Proj{\OK[X,Y,Z]/(F)}.
    \end{align}
    Is is clearly a proper $\OK$-scheme.
    For the flatness, it suffices to  show that $R \coloneqq \OK[X,Y,Z]/(F(X,Y,Z))$ is a torsion-free $\OK$-module by Remark \ref{rem:flatIffIntegral}.
    This is the case if and only if $\bar{F} \neq 0 \in k[X,Y,Z]$, which is equivalent to $F$ being primitive.
    For the normality, we use the following argument which is similar to Example 8.2.26 of \cite{Liu}.
    Since $\X$ is a Cohen-Macaulay scheme, it is normal if and only if it is normal at the points of codimension $1$.
    At this points, normality is equivalent to regularity.
    These are exactly the closed points on the generic fibre $\X_\eta=X$ and the generic points of the special fibre $\X_s$.
    Let $\xi \in \X$ be a generic point of the special fibre $\X_s$ and let $\pi\in \OK$ be a uniformising parameter.
    On the affine chart $U: Z \neq 0$ with $x = X/Z$ and $y=Y/Z$, the point $\xi$ corresponds to the ideal $\m$ that is generated by $\pi$ and a
    polynomial $G(x,y) \in \OK[x,y]$ whose image $\bar{G}(x,y) \in k[x,y]$ is irreducible.
    We can write
    \begin{align}
        F(x,y) = G(x,y)^r H_1(x,y) + \pi^e H_2(x,y),
    \end{align}
    for $H_i(x,y)\in \OK[x,y]$ and $r,e \geq 1$ and with $\bar{H}_1(x,y) \notin \bar{G}(x,y)k[x,y]$ and $\bar{H}_2(x,y) \neq 0$.
    By Corollary 4.2.12 of \cite{Liu}, $\xi$ is regular if $F \notin \m^2$.
    Thus, $\X$ is normal at $\xi$ if and only if either $r=1$; or $e=1$ and $\bar{H}_2(x,y) \neq 0$.
    \end{example}
\end{comment}

\begin{example}\label{ex:model-of-plane-curve}
    We discuss the important example of plane curves.
    A smooth and absolutely irreducible curve in $\PP^2_K$ is defined by a homogeneous irreducible polynomial $F \in K [X,Y,Z]$.
    After multiplying $F$ by a suitable element of $K^\times$, we may assume that $F$ is integral (i.e., all coefficients are in $\OK$) and primitive.
    Then the $\OK$-scheme
    \begin{align}
        \X = \Proj{\OK[X,Y,Z]/(F)}
    \end{align}
    is an obvious candidate for a model of $X$.
    $\X$ is proper as a closed subscheme of $\PP^2_K$ and flat by Remark \ref{rem:flatIffIntegral} since
    $\OK[X,Y,Z]/(F)$ is torsion-free by irreducibility and primitivity of $F$.
    Hence, $\X$ is a model if and only if it is normal.
    We rephrase the discussion of normality in \cite[Example 8.2.26]{Liu} for the present setting.
    Since $\X$ is Cohen--Macaulay, normality of $\X$ is equivalent to normality at the points of codimension $1$,
    which in turn is equivalent to regularity at these points.
    The points of codimension $1$ are the closed points of the (smooth) generic fibre $\X_\eta = X$,
    and the generic points of the special fibre $\X_s$.
    Now let $\xi \in \X_s$ be a generic point.
    On the affine chart $Z \neq 0$, with $x = X/Z$ and $y = Y/Z$, the point $\xi$ corresponds to the ideal $\m$ generated by a uniformising parameter $\pi \in \OK$
    and a polynomial $G \in K[x,y]$ whose reduction $\bar{G} \in k[x,y]$ is irreducible.
    We can write
    \begin{align}
        F(x,y) = G(x,y)^r H_1(x,y) + \pi^e H_2(x,y),
    \end{align}
    where $\bar{H}_1(x,y) \notin \bar{G}(x,y)k[x,y]$ and $\bar{H}_2(x,y)\neq 0$.
    By \cite[Corollary 4.2.12]{Liu}, $\X$ is regular at $\xi$ if and only if $F \notin \m^2$,
    which in turn is the case if and only if either $r=1$, or $e=1$ and $\bar{H}_2 \notin \bar{G}(x,y)k[x,y]$.

\end{example}

\begin{example}\label{ex:particular-model-of-plane-curve}
    \begin{enumerate}[label=\textup{(\alph*)}]
        \item Consider the curve $X$ defined by $y^2 = x^3 + px$ over $\QQ_p$ with $p \neq 2$.
            We have seen that $\X = \Proj{\OK[X,Y,Z]/(X^3 + pXZ^2 - Y^2Z)}$ is a proper and flat $\OK$-scheme with $\X_\eta = X$.
            The special fibre is given by $\X_s = \Proj{k[X,Y,Z]/(X^3 - Y^2Z)}$.
            On the affine cover $Z\neq 0$, its generic point $\xi$ corresponds to the ideal $\m = (p, x^3-y^2)$.
            Thus, in the notation of Example \ref{ex:model-of-plane-curve}, $r=e=1$; that is, $\X$ is normal and therefore a model of $X$.
        \item Let $X: y^3 = 1 + x^3 + 3x$ be a curve over $\QQ_3$.
            The scheme
            \begin{align}
                \X = \Proj{\OK[X,Y,Z]/(Y^3 - Z^3 - X^3 - 3XZ^2)}
            \end{align}
            has a non-reduced special fibre:
            On the affine chart $Z \neq 0$ we have $\X_s: 0 = y^3 - x^3 - 1 = (y - x - 1)^3$ over $\FF_3$.
            However, $\X$ is normal and thus a model of $X$ since
            \begin{align*}
                F(x,y) &= G(x,y)^r H_1(x,y) + 3 H_2(x,y) \\
                    &\coloneqq (y-x-1)^3 + 3(-x^2y + xy^2 + x^2 - 2xy + y^2 - y),
            \end{align*}
            i.e. $r=3$ but $e=1$ and $\bar{H}_2(x,y) \notin \bar{G}(x,y)\FF_3[x,y]$.
    \end{enumerate}
\end{example}

% Model with reduced special fibre: X: 0 = (y-1)(x^3+2) +3x

\section{Components and geometric valuations}\label{sec:components-and-geometric-valuations}
We now associate a valuation of the function field to each irreducible component of the special fibre. This construction also relates the multiplicity of a component to the value group of its valuation.

In this section, we assume the valuation $v_K$ to be \textit{normalized}, i.e. $v_K(K^\times) = \ZZ$.

\begin{definition}
    Let $F_X$ be the function field of $X$ over $K$.
    A valuation $v: F_X \to \hat{\RR}$ is said to be \textit{of type \textrm{II}}, if
    \begin{enumerate}[label=\textup{(\alph*)}]
        \item $v|_K = v_K$,
        \item $v$ is discrete, and
        \item the extension of residue fields $k(v)|k$ has transcendence degree $1$.
    \end{enumerate}
    A valuation of this type is also known as a \textit{geometric valuation}.
    We write $V_\textrm{II}(F_X)$ for the set of all valuations of type \textrm{II} of $F_X$.
    The number
    \begin{align}
        m(v) \coloneqq [\Gamma_v : \Gamma_K]
    \end{align}
    is called the \textit{multiplicity of $v$}, where $\Gamma_v \coloneqq v(F_X^*)$ is the value group.
    We say that $v$ is \textit{weakly unramified} if $m(v) = 1$.
\end{definition}

\begin{example}
    The \textit{Gauß valuation} on the polynomial ring $K[x]$ is defined as
    \begin{align}
        v_{\scriptscriptstyle G}: K[x] &\to \hat{\QQ}, \; v_{\scriptscriptstyle G}\left( \sum_i a_i x^{i} \right) = \min_i(v_K(a_i)).
    \end{align}
    Since $k(v_{\scriptscriptstyle G}) = k(\bar{x})$ (where $\bar{x}$ is the image of $x$ in $k(v_{\scriptscriptstyle G})$), the Gauß valuation
    extends naturally and uniquely to $K(x)$, so $v_{\scriptscriptstyle G} \in V_\textrm{II}(K(x))$.
\end{example}

Let $\X$ be a model of $X$.
We wish to assign a valuation in $V_\textrm{II}(F_X)$ to each generic point of the special fibre $\X_s$.
Suppose that $Z$ is an irreducible component of $\X_s$ with generic point $\xi$.
Since $\xi$ has codimension $1$ in $\X$, the stalk $\OO_{\X,\xi}$ is a discrete valuation ring\@.
Thus, we can define a valuation on $F_X$,
\begin{align} \label{eq:ValuationOfComponent}
    v_Z: \mathrm{Frac}(\OO_{\X, \xi}) = F_X \to \hat{\QQ}.
\end{align}
Since $\X$ is a flat $\OK $-scheme, $\OO_{\X, \xi}$ dominates $\OK$, and we may normalize $v_Z$ so that its restriction to $K$ is $v_K$.
We set
\begin{align}
    V(\X) \coloneqq \lbrace v_Z: \; Z \subset \X_s \text{ is an irreducible component}\rbrace.
\end{align}

\begin{proposition}\label{prop:WellDefOfInducedValuation}
    The valuation defined in~\eqref{eq:ValuationOfComponent} is of type \textrm{II}.
\end{proposition}

\begin{proof}
    As mentioned above, we can assume $v_Z|_K = v_K$.
    The residue field of the valuation is
    \begin{align}
        \kappa(v_Z)=\kappa(\xi)=k(Z).
    \end{align}
    Since $Z$ is an integral curve over $k$, it follows that
    $\trdeg_k \kappa(v_Z)=1$, showing $v_Z \in V_\textrm{II}(F_X)$.
\end{proof}

\begin{definition}
    The \textit{multiplicity} of the irreducible component $Z$ of $\X_s$ containing the generic point $\xi$ is defined as
    \begin{align}
        m(Z)
        \coloneqq \length_{\OO_{\X,\xi}}\bigl(\OO_{\X_s,\xi}\bigr)
        = \length_{\OO_{\X,\xi}}\bigl(\OO_{\X,\xi}/\pi\OO_{\X,\xi}\bigr),
    \end{align}
    where $\pi$ is a uniformising parameter of $\OK$.
\end{definition}

\begin{proposition}\label{prop:multiplicities-coincide}
    We have $m(v_Z) = m(Z)$.
\end{proposition}

\begin{proof}
    Since the model $\X$ is normal, the stalk $\OO_{\X,\xi}$ is a discrete valuation ring.
    Let $\pi_\xi$ be a uniformising parameter of this ring.
    There is a unit
    $u\in\OO_{\X,\xi}^{\times}$ such that
    \begin{align}
        \pi=u\pi_\xi^{m(Z)}.
    \end{align}
    Indeed, the exponent on the right is the length of
    $\OO_{\X,\xi}/\pi\OO_{\X,\xi}=\OO_{\X_s,\xi}$.
    Since $v_Z(\pi)=v_K(\pi)=1$, the displayed relation gives
    $v_Z(\pi_\xi)=1/m(Z)$.
    Hence
    \begin{align}
        m(Z)=[\Gamma_{v_Z}:\Gamma_K]=m(v_Z).
    \end{align}
\end{proof}

The finite set $V(\X)$ motivates the classification of models developed in the next chapter.
